\section{Challenge--Regulation--Recovery Event Model}

Let a recovery event $R_j$ be represented as
\[
R_j = \langle I_j, C_j, A_j, Q_j, O_j, V_j \rangle,
\]
where $I_j$ is the eligibility/reference state, $C_j$ the challenge description, $A_j$ the ordered regulation-action sequence, $Q_j$ the recovery-assessment trajectory, $O_j$ the terminal outcome and censoring state, and $V_j$ the version/provenance context.

Each element is represented by concrete telemetry events rather than inferred retrospectively from a single summary row. The reference implementation uses a JSON Schema contract and validates cross-field semantic invariants in addition to basic JSON types.

\subsection{Versioned structural inventory}
The recovery-event summary contract is distinct from the ordered telemetry stream and the controller replay fixtures. At the schema snapshot used for this manuscript, the top-level JSON Schema declares 41 properties, 32 of which are required, and the constant \texttt{schema\_version} is \texttt{0.6.0}. Table~\ref{tab:s02-schema-fields} is generated directly from the canonical schema (Git blob \texttt{0a7ee0a3c7e8db771d64c2aa188692fbe0f6c751}), not transcribed by hand. It summarizes only top-level declarations; nested constraints and cross-field semantic checks must be inspected in the schema and validator. This is a pinned descriptive snapshot for reproducibility, not a claim that the participant-facing protocol, export policy, or recovery construct has been prospectively frozen.

% Generated by scripts/generate_s02_schema_table.py; do not edit.
\begingroup
\footnotesize
\begin{longtable}{@{}p{.36\linewidth}p{.20\linewidth}p{.08\linewidth}p{.24\linewidth}@{}}
\caption{Top-level recovery-event schema fields (v0.6.0).}\label{tab:s02-schema-fields}\\
\toprule
Field & Type & Req. & Selected constraint \\
\midrule
\endfirsthead
\toprule
Field & Type & Req. & Selected constraint \\
\midrule
\endhead
\nolinkurl{schema_version} & string & yes & const: 0.6.0 \\
\nolinkurl{participant_pseudonym} & string & yes & min. length 1 \\
\nolinkurl{study_arm} & string & yes & enum (4) \\
\nolinkurl{intervention_version} & string & yes & min. length 1 \\
\nolinkurl{configuration_hash} & string / null & no & --- \\
\nolinkurl{session_id} & string & yes & min. length 1 \\
\nolinkurl{session_index} & integer & yes & min. 1 \\
\nolinkurl{task_id} & string & yes & min. length 1 \\
\nolinkurl{trial_id} & string / null & no & --- \\
\nolinkurl{challenge_event_id} & string & yes & min. length 1 \\
\nolinkurl{challenge_type} & string & yes & enum (5) \\
\nolinkurl{challenge_eligibility_rule_version} & string / null & no & --- \\
\nolinkurl{challenge_onset_ms} & integer & yes & min. 0 \\
\nolinkurl{reference_window} & object / null & no & nested object \\
\nolinkurl{cognitive_difficulty} & object & yes & --- \\
\nolinkurl{sensory_load} & object & yes & --- \\
\nolinkurl{cue_condition} & string / null & no & --- \\
\nolinkurl{explicit_child_input} & object / null & no & --- \\
\nolinkurl{recovery_parameters} & object & yes & nested object \\
\nolinkurl{recovery_measurement_context} & object & yes & nested object \\
\nolinkurl{recovery_measurement_context_unchanged} & boolean & yes & --- \\
\nolinkurl{recovery_data_valid} & boolean & yes & --- \\
\nolinkurl{invalid_reason} & string / null & yes & enum (7) \\
\nolinkurl{recovery_score_components} & object / null & yes & --- \\
\nolinkurl{recovery_observed} & boolean & yes & --- \\
\nolinkurl{recovery_latency_ms} & integer / null & yes & min. 0 \\
\nolinkurl{right_censored} & boolean & yes & --- \\
\nolinkurl{censor_reason} & string / null & no & enum (6) \\
\nolinkurl{telemetry_complete} & boolean & yes & --- \\
\nolinkurl{missingness_codes} & array & no & unique items \\
\nolinkurl{safety_override} & boolean & yes & --- \\
\nolinkurl{controller_reason_codes} & array & no & --- \\
\nolinkurl{recovery_feature_model_id} & string & yes & enum (4) \\
\nolinkurl{recovery_score_q} & number / null & yes & max. 0 \\
\nolinkurl{recovery_monitoring_onset_ms} & integer / null & yes & min. 0 \\
\nolinkurl{recovery_threshold_entry_ms} & integer / null & yes & min. 0 \\
\nolinkurl{recovery_confirmation_ms} & integer / null & yes & min. 0 \\
\nolinkurl{pre_monitoring_latency_ms} & integer / null & yes & min. 0 \\
\nolinkurl{total_return_latency_ms} & integer / null & yes & min. 0 \\
\nolinkurl{challenge_eligibility} & object & yes & nested object \\
\nolinkurl{target_rule_recovery_reference} & object & yes & nested object \\
\bottomrule
\end{longtable}
\endgroup


\subsection{Eligibility and reference state}
Before a challenge is eligible, the system records the task rule, immediate eligibility window, separate recovery-reference identity, timing-feature eligibility, and frozen challenge-rule version. This prevents a preswitch eligibility window from being silently reused as the recovery baseline when the protocol defines a different reference.

\subsection{Challenge and regulation}
Challenge onset receives a persistent event identifier. Regulation actions reference that identifier and include onset, action type, source of selection, and controller reason.

\subsection{Recovery trajectory}
Recovery is represented as repeated assessments of a versioned feature model rather than a single post hoc latency. The schema records component deficits, threshold, stability-window progress, and criterion state.

\subsection{Terminal outcome}
A terminal state specifies whether recovery was observed, right-censored, invalidated, interrupted, or unavailable because required telemetry was missing.

\subsection{Version context}
The event carries schema, intervention, configuration, device-qualification and measurement-rule versions so events remain interpretable after the application evolves.
